Fix \(D\in\mathcal D(G)\).  Let \(\mathbf M\in\mathcal C(\mathcal P)\) and put \(q=q_D^\star(\mathbf M)\).  The complete response-profile definition gives \(q\geq0\) and \(\mathbf 1^{\mathsf T}q=1\).  Each retained raw row has district-locality and transport-admissibility proofs, so typed-menu realizability, total probability, and district-kernel invariance give
\[
P_a^s(v)=(r^{\mathrm{raw}}_{s,a,D,v})^{\mathsf T}q.
\tag{1}
\]
For every retained \(\operatorname{ROW\_ELIGIBLE}\) cell, district-cell linearity gives the two separate equations
\[
N_{h,c}=u_{h,c}^{\mathsf T}q,
\qquad
H_h=v_h^{\mathsf T}q.
\tag{2}
\]
The conditional quotient is not used as a row.  By the stacking definition, (1), (2), and normalization are exactly the equations in \(A_Dq=m_D\).  Therefore
\[
\mathcal Q_D(\mathcal P)\subseteq\mathcal K_D(m_D).
\tag{3}
\]

Assume \(\mathsf{BASep}_D(\mathcal P)\).  Its definition supplies a single baseline \(\mathbf M^0\in\mathcal C_{\mathrm{sat},D}(\mathcal P)\).  Put
\[
K=\mathcal K_D(m_D),
\qquad q_0=q_D^\star(\mathbf M^0),
\qquad Z=Z_D(m_D).
\]
By (3), \(q_0\in K\), so \(K\ne\varnothing\).

Let \(J=\Omega_D\setminus Z\).  For every \(\omega\in J\), the definition of \(Z\) supplies \(q^\omega\in K\) with \(q^\omega(\omega)>0\).  Hence
\[
\bar q=|J|^{-1}\sum_{\omega\in J}q^\omega
\]
belongs to \(K\), vanishes on \(Z\), and is strictly positive on \(J\).  Define
\[
L=\{v:A_Dv=0,\ v(\omega)=0\text{ for }\omega\in Z\}.
\]
Every difference of two points of \(K\) lies in \(L\).  Conversely, for \(v\in L\), finiteness of \(J\) and strict positivity of \(\bar q\) on \(J\) permit \(\varepsilon>0\) with \(\bar q\pm\varepsilon v\geq0\).  Both perturbations lie in \(K\), and
\[
v=(2\varepsilon)^{-1}\{(\bar q+\varepsilon v)-(\bar q-\varepsilon v)\}.
\]
Thus
\[
L=\operatorname{span}(K-K).
\tag{4}
\]
Appending \(e_\omega^{\mathsf T}\), \(\omega\in Z\), to \(A_D\) gives a matrix \(\widetilde A_D\) with \(L=\ker\widetilde A_D\).  Finite-dimensional row reduction therefore yields
\[
R_D(m_D)=\operatorname{row}(\widetilde A_D)
=(\ker\widetilde A_D)^\perp
=\operatorname{span}(K-K)^\perp.
\tag{5}
\]

Fix arbitrary \(q\in K\).  Since \(\Omega_D=\Omega_D^{\mathrm{all}}\), replace the target law of the full focal block \(U_D^\star\) by \(q\), and make the same replacement in every source satisfying \(D\cap\Delta_s=\varnothing\).  Keep the baseline focal law in every source satisfying \(D\subseteq\Delta_s\), and keep every nonfocal mechanism and latent law fixed.  District closure makes these source cases exhaustive.  The unrestricted focal-block clause permits the replacement on the original observed alphabets.  Independence of the focal block from nonfocal blocks preserves district factorization; common replacement in the target and unchanged sources preserves district-kernel invariance.  Call the resulting family \(\mathbf M^q\).

Consider a transport-admissible complete supplied cell \(C_{a,v}\) in domain \(s\).  The fixed nonfocal product law is \(\nu_{-D}^{0,s}\), and finite total probability gives
\[
\begin{aligned}
P_{M^{q,s}}^a(v)-P_{M^{0,s}}^a(v)
&=\sum_{\omega_{-D}}\nu_{-D}^{0,s}(\omega_{-D})
\chi_{C_{a,v},\omega_{-D}}^{\mathsf T}(q-q_0)\\
&=(\bar\chi_{C_{a,v}}^{s,\mathbf M^0})^{\mathsf T}(q-q_0).
\end{aligned}
\tag{6}
\]
The witnessing baseline-averaged condition puts the row in \(R_D(m_D)\), while (4) puts \(q-q_0\) in its orthogonal complement by (5).  Thus (6) is zero.  A non-admissible complete cell occurs in a changed source; its focal and nonfocal laws were both left fixed, so it too is preserved.  Hence all complete supplied laws remain equal to \(\mathcal P\), and
\[
\mathbf M^q\in\mathcal C_{\mathrm{sat},D}(\mathcal P),
\qquad q_D^\star(\mathbf M^q)=q.
\]
Since \(q\in K\) was arbitrary, (3) gives
\[
\mathcal Q_D^{\mathrm{sat},D}(\mathcal P)=K.
\]
The construction includes every boundary point, changes only the focal district, and retains every nonfocal graph-compatible restriction.  The established implication from active-slice to baseline-averaged separation gives the sufficient certificate consisting of nonemptiness plus \(\mathsf{ASep}_D\); the established structural-separation implication gives the further sufficient \(\mathsf{Sep}_D\) check.